% Manual tecnico del proyecto final de Computacion Grafica
% Se compila con latexmk (ver latexmkrc). Las capturas se toman de ../docs/evidencias/ y de img/
% Lo que falta se marca con \pendiente{...}
\documentclass[11pt,a4paper]{article}
% Estilo comun de los manuales, el mismo de la propuesta y del reporte de avances:
% fondo oscuro, titulos morados y recuadro de encabezado. Se carga con % Estilo comun de los manuales, el mismo de la propuesta y del reporte de avances:
% fondo oscuro, titulos morados y recuadro de encabezado. Se carga con % Estilo comun de los manuales, el mismo de la propuesta y del reporte de avances:
% fondo oscuro, titulos morados y recuadro de encabezado. Se carga con \input{estilo}.
\usepackage[utf8]{inputenc}
\usepackage[T1]{fontenc}
\usepackage[spanish,es-noshorthands,es-tabla]{babel}
\usepackage{lmodern}
\usepackage[margin=2.5cm]{geometry}
\usepackage{graphicx}
\usepackage{subcaption}
\usepackage{float}
\usepackage[table]{xcolor}
\usepackage{titlesec}
\usepackage{tcolorbox}
\usepackage{enumitem}
\usepackage{array}
\usepackage{longtable}
\usepackage{amssymb}
\usepackage{pgfgantt}
\usepackage{hyperref}

\definecolor{fondo}{HTML}{1E1E2E}
\definecolor{texto}{HTML}{CDCDD6}
\definecolor{morado}{HTML}{B48EDB}
\definecolor{azul}{HTML}{8AADF4}
\definecolor{rosa}{HTML}{E3A8D6}
\definecolor{verde}{HTML}{A6DA95}
\definecolor{linea}{HTML}{55556A}

\pagecolor{fondo}
\color{texto}
\hypersetup{colorlinks=true, linkcolor=azul, urlcolor=azul, citecolor=azul}
\arrayrulecolor{texto}
\graphicspath{{../docs/evidencias/}{img/}}

\setlength{\parindent}{0pt}
\setlength{\parskip}{6pt}

% secciones numeradas, moradas y con una linea abajo
\titleformat{\section}{\Large\bfseries\color{morado}}{\thesection.}{8pt}{}[{\color{linea}\titlerule}]
\titleformat{\subsection}{\large\bfseries\color{azul}}{\thesubsection}{8pt}{}
\titleformat{\subsubsection}{\normalsize\bfseries\color{rosa}}{}{0pt}{}
\titlespacing*{\section}{0pt}{18pt}{10pt}
\titlespacing*{\subsection}{0pt}{12pt}{6pt}

\setlist[itemize,1]{label=\scriptsize$\blacksquare$}
\setlist{itemsep=3pt}

% lo que falta por escribir o por confirmar; se busca con "pendiente" para no dejar ninguno
\newcommand{\pendiente}[1]{%
  \begin{tcolorbox}[colback=fondo, colframe=rosa, coltext=rosa, arc=2mm, boxrule=0.8pt, left=6pt, right=6pt, top=4pt, bottom=4pt]
    \textbf{Pendiente:} #1
  \end{tcolorbox}}

% tecla con recuadro
\newcommand{\tecla}[1]{\tcbox[on line, colback=fondo, colframe=azul, coltext=azul, arc=1mm, boxrule=0.6pt, left=2pt, right=2pt, top=0pt, bottom=0pt]{\small\ttfamily #1}}

\newcommand{\dosfotos}[5]{%
  \begin{figure}[H]
    \centering
    \begin{subfigure}{0.48\textwidth}
      \includegraphics[width=\linewidth]{#1}
      \caption{#2}
    \end{subfigure}\hfill
    \begin{subfigure}{0.48\textwidth}
      \includegraphics[width=\linewidth]{#3}
      \caption{#4}
    \end{subfigure}
    \caption{#5}
  \end{figure}}

\newcommand{\unafoto}[2]{%
  \begin{figure}[H]
    \centering
    \includegraphics[width=0.8\textwidth]{#1}
    \caption{#2}
  \end{figure}}

% encabezado en recuadro: titulo del documento
\newcommand{\encabezado}[1]{%
  \begin{tcolorbox}[colback=fondo, colframe=morado, coltext=texto, arc=4mm, boxrule=0.8pt, left=10pt, right=10pt, top=8pt, bottom=8pt]
    \centering
    {\Large\bfseries\color{azul} #1 $\|$ Calenda de Oaxaca}\\[4pt]
    {\large\bfseries Facultad de Ingeniería -- UNAM}\\[10pt]
    \begin{tabular}{@{}p{0.47\linewidth}p{0.47\linewidth}@{}}
      {\color{rosa}\textbf{Equipo:}} 16 & {\color{rosa}\textbf{Materia:}} Computación Gráfica e I.H.C. \\
      {\color{rosa}\textbf{Número de cuenta:}} 421112601 & {\color{rosa}\textbf{Grupo:}} 01 \\
      {\color{rosa}\textbf{Semestre:}} 2027-1 & {\color{rosa}\textbf{Entrega:}} 18 de noviembre de 2026 \\
    \end{tabular}
  \end{tcolorbox}}
.
\usepackage[utf8]{inputenc}
\usepackage[T1]{fontenc}
\usepackage[spanish,es-noshorthands,es-tabla]{babel}
\usepackage{lmodern}
\usepackage[margin=2.5cm]{geometry}
\usepackage{graphicx}
\usepackage{subcaption}
\usepackage{float}
\usepackage[table]{xcolor}
\usepackage{titlesec}
\usepackage{tcolorbox}
\usepackage{enumitem}
\usepackage{array}
\usepackage{longtable}
\usepackage{amssymb}
\usepackage{pgfgantt}
\usepackage{hyperref}

\definecolor{fondo}{HTML}{1E1E2E}
\definecolor{texto}{HTML}{CDCDD6}
\definecolor{morado}{HTML}{B48EDB}
\definecolor{azul}{HTML}{8AADF4}
\definecolor{rosa}{HTML}{E3A8D6}
\definecolor{verde}{HTML}{A6DA95}
\definecolor{linea}{HTML}{55556A}

\pagecolor{fondo}
\color{texto}
\hypersetup{colorlinks=true, linkcolor=azul, urlcolor=azul, citecolor=azul}
\arrayrulecolor{texto}
\graphicspath{{../docs/evidencias/}{img/}}

\setlength{\parindent}{0pt}
\setlength{\parskip}{6pt}

% secciones numeradas, moradas y con una linea abajo
\titleformat{\section}{\Large\bfseries\color{morado}}{\thesection.}{8pt}{}[{\color{linea}\titlerule}]
\titleformat{\subsection}{\large\bfseries\color{azul}}{\thesubsection}{8pt}{}
\titleformat{\subsubsection}{\normalsize\bfseries\color{rosa}}{}{0pt}{}
\titlespacing*{\section}{0pt}{18pt}{10pt}
\titlespacing*{\subsection}{0pt}{12pt}{6pt}

\setlist[itemize,1]{label=\scriptsize$\blacksquare$}
\setlist{itemsep=3pt}

% lo que falta por escribir o por confirmar; se busca con "pendiente" para no dejar ninguno
\newcommand{\pendiente}[1]{%
  \begin{tcolorbox}[colback=fondo, colframe=rosa, coltext=rosa, arc=2mm, boxrule=0.8pt, left=6pt, right=6pt, top=4pt, bottom=4pt]
    \textbf{Pendiente:} #1
  \end{tcolorbox}}

% tecla con recuadro
\newcommand{\tecla}[1]{\tcbox[on line, colback=fondo, colframe=azul, coltext=azul, arc=1mm, boxrule=0.6pt, left=2pt, right=2pt, top=0pt, bottom=0pt]{\small\ttfamily #1}}

\newcommand{\dosfotos}[5]{%
  \begin{figure}[H]
    \centering
    \begin{subfigure}{0.48\textwidth}
      \includegraphics[width=\linewidth]{#1}
      \caption{#2}
    \end{subfigure}\hfill
    \begin{subfigure}{0.48\textwidth}
      \includegraphics[width=\linewidth]{#3}
      \caption{#4}
    \end{subfigure}
    \caption{#5}
  \end{figure}}

\newcommand{\unafoto}[2]{%
  \begin{figure}[H]
    \centering
    \includegraphics[width=0.8\textwidth]{#1}
    \caption{#2}
  \end{figure}}

% encabezado en recuadro: titulo del documento
\newcommand{\encabezado}[1]{%
  \begin{tcolorbox}[colback=fondo, colframe=morado, coltext=texto, arc=4mm, boxrule=0.8pt, left=10pt, right=10pt, top=8pt, bottom=8pt]
    \centering
    {\Large\bfseries\color{azul} #1 $\|$ Calenda de Oaxaca}\\[4pt]
    {\large\bfseries Facultad de Ingeniería -- UNAM}\\[10pt]
    \begin{tabular}{@{}p{0.47\linewidth}p{0.47\linewidth}@{}}
      {\color{rosa}\textbf{Equipo:}} 16 & {\color{rosa}\textbf{Materia:}} Computación Gráfica e I.H.C. \\
      {\color{rosa}\textbf{Número de cuenta:}} 421112601 & {\color{rosa}\textbf{Grupo:}} 01 \\
      {\color{rosa}\textbf{Semestre:}} 2027-1 & {\color{rosa}\textbf{Entrega:}} 18 de noviembre de 2026 \\
    \end{tabular}
  \end{tcolorbox}}
.
\usepackage[utf8]{inputenc}
\usepackage[T1]{fontenc}
\usepackage[spanish,es-noshorthands,es-tabla]{babel}
\usepackage{lmodern}
\usepackage[margin=2.5cm]{geometry}
\usepackage{graphicx}
\usepackage{subcaption}
\usepackage{float}
\usepackage[table]{xcolor}
\usepackage{titlesec}
\usepackage{tcolorbox}
\usepackage{enumitem}
\usepackage{array}
\usepackage{longtable}
\usepackage{amssymb}
\usepackage{pgfgantt}
\usepackage{hyperref}

\definecolor{fondo}{HTML}{1E1E2E}
\definecolor{texto}{HTML}{CDCDD6}
\definecolor{morado}{HTML}{B48EDB}
\definecolor{azul}{HTML}{8AADF4}
\definecolor{rosa}{HTML}{E3A8D6}
\definecolor{verde}{HTML}{A6DA95}
\definecolor{linea}{HTML}{55556A}

\pagecolor{fondo}
\color{texto}
\hypersetup{colorlinks=true, linkcolor=azul, urlcolor=azul, citecolor=azul}
\arrayrulecolor{texto}
\graphicspath{{../docs/evidencias/}{img/}}

\setlength{\parindent}{0pt}
\setlength{\parskip}{6pt}

% secciones numeradas, moradas y con una linea abajo
\titleformat{\section}{\Large\bfseries\color{morado}}{\thesection.}{8pt}{}[{\color{linea}\titlerule}]
\titleformat{\subsection}{\large\bfseries\color{azul}}{\thesubsection}{8pt}{}
\titleformat{\subsubsection}{\normalsize\bfseries\color{rosa}}{}{0pt}{}
\titlespacing*{\section}{0pt}{18pt}{10pt}
\titlespacing*{\subsection}{0pt}{12pt}{6pt}

\setlist[itemize,1]{label=\scriptsize$\blacksquare$}
\setlist{itemsep=3pt}

% lo que falta por escribir o por confirmar; se busca con "pendiente" para no dejar ninguno
\newcommand{\pendiente}[1]{%
  \begin{tcolorbox}[colback=fondo, colframe=rosa, coltext=rosa, arc=2mm, boxrule=0.8pt, left=6pt, right=6pt, top=4pt, bottom=4pt]
    \textbf{Pendiente:} #1
  \end{tcolorbox}}

% tecla con recuadro
\newcommand{\tecla}[1]{\tcbox[on line, colback=fondo, colframe=azul, coltext=azul, arc=1mm, boxrule=0.6pt, left=2pt, right=2pt, top=0pt, bottom=0pt]{\small\ttfamily #1}}

\newcommand{\dosfotos}[5]{%
  \begin{figure}[H]
    \centering
    \begin{subfigure}{0.48\textwidth}
      \includegraphics[width=\linewidth]{#1}
      \caption{#2}
    \end{subfigure}\hfill
    \begin{subfigure}{0.48\textwidth}
      \includegraphics[width=\linewidth]{#3}
      \caption{#4}
    \end{subfigure}
    \caption{#5}
  \end{figure}}

\newcommand{\unafoto}[2]{%
  \begin{figure}[H]
    \centering
    \includegraphics[width=0.8\textwidth]{#1}
    \caption{#2}
  \end{figure}}

% encabezado en recuadro: titulo del documento
\newcommand{\encabezado}[1]{%
  \begin{tcolorbox}[colback=fondo, colframe=morado, coltext=texto, arc=4mm, boxrule=0.8pt, left=10pt, right=10pt, top=8pt, bottom=8pt]
    \centering
    {\Large\bfseries\color{azul} #1 $\|$ Calenda de Oaxaca}\\[4pt]
    {\large\bfseries Facultad de Ingeniería -- UNAM}\\[10pt]
    \begin{tabular}{@{}p{0.47\linewidth}p{0.47\linewidth}@{}}
      {\color{rosa}\textbf{Equipo:}} 16 & {\color{rosa}\textbf{Materia:}} Computación Gráfica e I.H.C. \\
      {\color{rosa}\textbf{Número de cuenta:}} 421112601 & {\color{rosa}\textbf{Grupo:}} 01 \\
      {\color{rosa}\textbf{Semestre:}} 2027-1 & {\color{rosa}\textbf{Entrega:}} 18 de noviembre de 2026 \\
    \end{tabular}
  \end{tcolorbox}}


\begin{document}

\encabezado{Manual técnico}

\tableofcontents

% =====================================================================
\section{Descripción general}

El proyecto es un recorrido virtual en 3D por el Centro Histórico de Oaxaca durante una calenda, hecho en C++ con OpenGL.
Una calenda es un desfile con banda de música, chinas oaxaqueñas, marmotas y toritos que va por las calles anunciando una fiesta. Hay calendas todo el año, pero la más conocida es la de la Guelaguetza, en julio.

En la escena, el desfile sigue el mismo camino que hacen muchas calendas en la ciudad. Sale del atrio del templo de Santo Domingo de Guzmán, baja por el andador Macedonio Alcalá pasando frente al templo de la Sangre de Cristo y el jardín Labastida, y termina en el Zócalo, donde se detiene frente al kiosco.
Mientras tanto, el usuario puede caminar por donde quiera, cambiar entre cinco cámaras, pasar de la tarde a la noche, seguir al desfile y buscar las cartas que tienen información de los edificios, los personajes y las artesanías.

\unafoto{retoque_texturas/05_aerea}{La escena completa vista desde arriba: el atrio de Santo Domingo arriba a la derecha, el andador al centro y el Zócalo abajo.}

\subsection{Cambios respecto a la propuesta}

\begin{itemize}
  \item \textbf{Recorrido.} Al principio la propuesta era hacer un circuito cerrado frente a Santo Domingo, pero al comparar la escena con fotos y mapas del lugar se decidió cambiarlo por algo más parecido al andador Macedonio Alcalá, que es por donde suelen pasar las calendas que salen de Santo Domingo hacia el Zócalo. El cambio se hizo el 2 de noviembre, y desde entonces el desfile ya no da vueltas, sino que va en una sola dirección, del atrio al Zócalo.
  \item \textbf{Cámaras.} En la propuesta solo había tres cámaras (la libre, la de la fachada y la del desfile). Al final se agregaron dos más, una en el Zócalo y otra aérea, y también una tecla para soltar cualquier cámara fija y seguir caminando desde donde estaba.
  \item \textbf{Animaciones.} En la propuesta el torito giraba sobre el danzante, pero se cambió por algo más parecido a uno real, donde el danzante lo sujeta de los carrizos de los lados mientras lo mece de lado a lado, y el torito tiene dos ruedas de cohetes que van dando vueltas y echando chispas. También se decidió que los personajes no movieran los brazos al caminar, porque todos van cargando algo. Además de la china, la marmota también se animó con KeyFrames, y se agregaron teclas para pausar y reiniciar el desfile.
  \item \textbf{Información cultural.} En la propuesta la información se iba a mostrar en tarjetas que se abrían con una tecla, pero se cambió por cartas de lotería repartidas por la escena que hay que ir encontrando. Al tomar una se puede voltear para leer su ficha mientras se escucha su narración, y las cartas de los personajes del desfile van caminando junto a ellos.
  \item \textbf{Audio.} Inicialmente solo se tenía contemplada la banda de viento sonando en bucle, pero se agregó audio espacial para que la banda suene más fuerte o más bajo según la cámara se acerca o se aleja del músico de en medio, que es de donde sale la música. También se integró un murmullo de gente que se repite en bucle mientras corre el programa, y se añadieron efectos de sonido, como las chispas del torito, las campanas y un cohete cuando el desfile sale de Santo Domingo, y otro cohete cuando llega al kiosco. Por último, se agregó una tecla (\tecla{M}) para quitar el sonido.
  \item \textbf{Velocidad del desfile.} Se subió la velocidad del desfile a 2 metros por segundo, para que la calenda tarde unos dos minutos en llegar al Zócalo.
  \item \textbf{Modo Día de Muertos.} Esto no estaba en la propuesta y se agregó al final como un extra. Con la tecla \tecla{O}, la misma escena se adorna como si fuera una noche de Día de Muertos en Oaxaca (sección \ref{sec:modomuertos}).
\end{itemize}

% =====================================================================
\section{Investigación cultural}
\label{sec:cultura}

\subsection{La Guelaguetza y la calenda}

La Guelaguetza es la fiesta más importante de Oaxaca. Su nombre viene del zapoteco y se suele traducir como ofrenda, cooperación o intercambio entre las comunidades (Embajada de México ante la Santa Sede, s.~f.). Se celebra los dos lunes que siguen a la fiesta de la Virgen del Carmen (16 de julio), y por eso también se le llama Lunes del Cerro. Esos días, las delegaciones de las regiones del estado bailan en el auditorio del Cerro del Fortín y al terminar le regalan al público productos de su región (Gobierno del Estado de Oaxaca, 2026).

Antes de cada fiesta se hace una calenda, que es un desfile por las calles para anunciar la celebración e invitar a la gente a unirse. En Oaxaca hay calendas casi todos los días, ya sea por la fiesta de un barrio, una boda o un bautizo, y en la Guelaguetza las delegaciones desfilan por el centro de la ciudad (México Desconocido, s.~f.). Es una costumbre tan común que hasta las instituciones la usan para festejar. Por ejemplo, en 2023 el INAH celebró los 50 años de su centro en Oaxaca con una calenda de más de 450 personas que recorrió las calles principales y terminó en la plaza de Santo Domingo con fuegos artificiales y un castillo (Instituto Nacional de Antropología e Historia [INAH], 2023). En el programa también truenan cohetes cuando el desfile sale y cuando llega al kiosco. En casi todas las calendas aparecen estos elementos:

\begin{itemize}
  \item \textbf{Banda de música de viento}, que toca sones y jarabes y va marcando el paso del desfile.
  \item \textbf{Chinas oaxaqueñas}, mujeres con el traje tradicional (blusa bordada, enagua de colores con olanes y rebozo) que bailan con una canasta adornada con flores en la cabeza y a veces reparten dulces o mezcal.
  \item \textbf{Marmotas}, esferas grandes de manta sobre un armazón de carrizo, montadas en un palo que un danzante va girando mientras avanza. Muchas llevan escrito el nombre del barrio o de la fiesta.
  \item \textbf{Toritos}, armazones de carrizo con forma de toro que un danzante carga sobre los hombros y que en la noche llevan cohetes.
  \item \textbf{Monos de calenda}, figuras gigantes de carrizo y tela que bailan girando.
  \item \textbf{Espectadores} a los lados de la calle, que aplauden, saludan y a veces se unen al desfile.
\end{itemize}

\subsection{Por qué se eligieron estos lugares}

\begin{itemize}
  \item \textbf{Templo de Santo Domingo de Guzmán.} Lo construyeron los dominicos entre los siglos XVI y XVIII, y su ex convento hoy es el Museo de las Culturas de Oaxaca (INAH, s.~f.). De su atrio salen muchas calendas, por eso ahí empieza el recorrido. La fachada de cantera verde, las dos torres y las jardineras con agaves del atrio se hicieron a partir de fotos del lugar.
  \item \textbf{Andador Macedonio Alcalá.} Es la calle peatonal que va de Santo Domingo al Zócalo y por donde pasan los desfiles. Tiene casas de colores de uno y dos pisos con balcones de herrería, y el piso es de cantera. Su nombre es por Macedonio Alcalá, el compositor oaxaqueño de ``Dios nunca muere''.
  \item \textbf{Templo de la Sangre de Cristo y jardín Labastida.} El templo está sobre el andador y el jardín a media cuadra, así que los dos quedan en el camino del desfile. En el jardín suelen ponerse artesanos y vendedores.
  \item \textbf{Zócalo.} Es la plaza principal de la ciudad. Tiene el kiosco en medio, árboles grandes, bancas de hierro y portales alrededor, y es donde termina el recorrido.
\end{itemize}

El Centro Histórico de Oaxaca es Patrimonio de la Humanidad desde 1987, junto con la zona arqueológica de Monte Albán (UNESCO World Heritage Centre, s.~f.).

\subsection{Objetos y adornos}

\begin{itemize}
  \item \textbf{Papel picado}, que se cuelga sobre las calles en las fiestas. En el programa son tiras de banderitas de colores con los recortes transparentes.
  \item \textbf{Artesanías}. Se pusieron alebrijes, que son figuras de madera de copal talladas y pintadas, como las que se hacen en talleres familiares de San Martín Tilcajete (Secretaría de Turismo de Oaxaca, 2018a), y piezas de barro. También hay tapetes de lana de Teotitlán del Valle, que se tejen en telar de pedales, se tiñen con tintes naturales como la grana cochinilla y llevan grecas zapotecas. Desde 2022 estos tapetes tienen una indicación geográfica protegida (Instituto Mexicano de la Propiedad Industrial [IMPI], 2022; Secretaría de Turismo de Oaxaca, 2018b). Todo está en puestos sobre petates y en mesas con toldo, como los que se ven en el andador y en el jardín Labastida.
  \item \textbf{Comida}: pan de yema en canastas y carritos de elotes en el Zócalo.
  \item \textbf{Plantas}: agaves en el atrio de Santo Domingo, flamboyanes de flores rojas, jacarandas, bugambilias y cacaloxóchitl (flor de mayo), que son comunes en el centro de Oaxaca.
\end{itemize}

\subsubsection{Época del año}

Aunque hay calendas todo el año, la más conocida es la de julio, así que se revisó si lo que aparece en la escena tendría sentido en esa época:

\begin{itemize}
  \item El pan de yema se vende todo el año, pero el pan de muerto con carita es de noviembre, por eso solo aparece en el modo Día de Muertos.
  \item Los elotes son de la temporada de lluvias (de julio a septiembre), justo cuando es la Guelaguetza.
  \item Los flamboyanes florecen al empezar las lluvias, entre mayo y julio, y las bugambilias tienen flor casi todo el año.
  \item Las jacarandas florecen en primavera, de febrero a abril, y en julio ya solo tienen hojas. Aun así se dejaron con flores porque son de los árboles más reconocibles del centro de Oaxaca y del jardín Labastida, aunque en una calenda de julio estarían verdes.
\end{itemize}

\subsection{El Día de Muertos en Oaxaca}

El modo extra muestra la misma calle en las noches del 1 y 2 de noviembre. Desde 2008, las fiestas indígenas dedicadas a los muertos están en la Lista Representativa del Patrimonio Cultural Inmaterial de la UNESCO. Según la tradición, en esos días las almas de los difuntos regresan a visitar a su familia, y para recibirlas se ponen ofrendas con flores, velas, comida y fotos, y se marca el camino con pétalos de cempasúchil (UNESCO, 2008). En Oaxaca estas fechas también se viven en la calle, con comparsas y calendas de disfraces. Con eso en mente se pusieron estos elementos:

\begin{itemize}
  \item \textbf{Cempasúchil} en cortinas sobre el andador, en esferas y guirnaldas en el kiosco y en un camino de pétalos que lleva al altar.
  \item \textbf{Velas}: veladoras en las ofrendas y velas flotantes colgadas entre el papel picado, que cambia por uno con calaveras, velas y cruces.
  \item \textbf{Un altar escalonado} junto al kiosco, con grecas como las de Mitla, arco de flores, macetas de cempasúchil, calaveritas de azúcar y un tapete de arena como los que se hacen en las calles de Oaxaca.
  \item \textbf{Ofrendas familiares} con mantel, veladoras, pan de muerto y calaveritas.
  \item \textbf{La catrina} y \textbf{el xoloitzcuintle}, que según la tradición es el perro que acompaña a las almas.
  \item \textbf{Caras de calavera} en los espectadores y en los personajes del desfile, como en las comparsas.
\end{itemize}

\subsection{La lotería y la información cultural}
\label{sec:loteria}

Para mostrar la información del recorrido se usaron cartas de la lotería mexicana. La lotería es un juego de figuras que se hizo popular en México a finales del siglo XIX, y la baraja más conocida, de 54 figuras, se le atribuye al francés Clemente Jacques, que la imprimió y la difundió en el centro del país (Sánchez Ulloa, 2022). Desde entonces se juega en las casas y en los puestos de las ferias y de las fiestas de barrio, así que mucha gente la relaciona con la fiesta. También se eligió porque la idea de llenar el tablero hace que el usuario quiera encontrar todas las cartas, y para eso tiene que recorrer toda la escena.

Cada carta se escogió por algo que la relaciona con el tema que explica:

\begin{table}[H]
  \centering
  \small
  \begin{tabular}{|p{0.17\linewidth}|p{0.28\linewidth}|p{0.45\linewidth}|}
    \hline
    \textbf{Carta} & \textbf{Tema y dónde está} & \textbf{Por qué esa carta} \\ \hline
    El Barril & La calenda, en el atrio & En la calenda se reparte mezcal \\ \hline
    La Luna & La marmota, con su danzante & Es redonda y lleva una luz por dentro \\ \hline
    La Dama & La china oaxaqueña, con la primera china & Representa a la mujer con su traje de fiesta \\ \hline
    El Diablito & El torito, con su danzante & Por los cohetes y la lumbre \\ \hline
    El Tambor & La banda de viento, con el músico de en medio & Es uno de los instrumentos de la banda \\ \hline
    El Sol & Santo Domingo, frente a la puerta & El interior del templo está cubierto de hoja de oro \\ \hline
    El Corazón & La Sangre de Cristo, en la banqueta & Por el nombre del templo \\ \hline
    El Catrín & El Zócalo, frente al kiosco & Por el paseo de domingo en la plaza \\ \hline
    La Mano & Los alebrijes, sobre su petate & Se tallan y se pintan a mano \\ \hline
    El Nopal & Los tapetes, junto a los textiles & La grana cochinilla vive en el nopal y da los rojos de los tapetes \\ \hline
  \end{tabular}
  \caption{Las diez cartas, su tema y dónde están.}
\end{table}

Los textos de las fichas son cortos, de tres o cuatro frases, para que se alcancen a leer mientras suena la narración, que se grabó con la voz del autor del proyecto. La información de las fichas se sacó de las fuentes de esta sección y de las páginas de los propios lugares: la historia de Santo Domingo y del Museo de las Culturas (INAH, s.~f.), el centro histórico y su traza (UNESCO World Heritage Centre, s.~f.), los alebrijes y los tapetes (IMPI, 2022; Secretaría de Turismo de Oaxaca, 2018a, 2018b) y la calenda y sus personajes (INAH, 2023; México Desconocido, s.~f.).

% =====================================================================
\section{Arquitectura del proyecto}

\subsection{Organización de carpetas}

\begin{table}[H]
  \centering
  \small
  \begin{tabular}{|>{\ttfamily}p{0.22\linewidth}|p{0.7\linewidth}|}
    \hline
    \textbf{Carpeta} & \textbf{Qué contiene} \\ \hline
    src/ & Todo el código en C++: la escena (\texttt{main.cpp}) y las clases de la ventana, la cámara, las luces, las texturas y los modelos \\ \hline
    shaders/ & Shaders en GLSL: \texttt{shader\_light} (escena con iluminación) y \texttt{skybox} \\ \hline
    Models/ & Modelos \texttt{.obj} con su \texttt{.mtl} (84 modelos, contando las piezas de los personajes y los del modo Día de Muertos) \\ \hline
    Textures/ & Texturas \texttt{.jpg} y \texttt{.png} (142 imágenes), las 18 caras de los tres cielos en \texttt{Skybox/} y las 35 imágenes de las cartas en \texttt{cartas/} \\ \hline
    Audio/ & Música de la banda, efectos (cohetes, campanas, chispas, murmullo) y en \texttt{narracion/} la narración de cada carta \\ \hline
    Librerias/ & GLEW, GLFW, GLM, Assimp, stb\_image y miniaudio (se reparte en \texttt{Librerias.zip}) \\ \hline
    docs/evidencias/ & Capturas de cada fase del desarrollo \\ \hline
    Documentacion/ & Estos manuales en LaTeX \\ \hline
  \end{tabular}
  \caption{Carpetas del proyecto.}
\end{table}

\subsection{Clases}

El código parte del que se usó en el laboratorio de la materia y conserva sus mismas clases. A algunas se les agregaron funciones cuando hizo falta, pero los modelos y las texturas se siguen cargando y dibujando igual que en el laboratorio.

\begin{table}[H]
  \centering
  \small
  \begin{tabular}{|>{\ttfamily}p{0.25\linewidth}|p{0.67\linewidth}|}
    \hline
    \textbf{Clase o archivo} & \textbf{Para qué sirve} \\ \hline
    main.cpp & Crea las mallas, carga shaders, texturas y modelos, coloca las luces y tiene el ciclo que dibuja cada cuadro \\ \hline
    Window & Crea la ventana con GLFW, lee el teclado y el mouse, y guarda la hora (tarde o noche), si está el modo Día de Muertos, la cámara activa y si el desfile está en pausa o se tiene que reiniciar \\ \hline
    Camera & Cámara en primera persona con WASD y mouse; \texttt{tomarVistaDe} copia la vista de una cámara fija y \texttt{limitarPosicion} evita que salga de la ciudad \\ \hline
    MeshModel (Mesh\_tn) & Mallas hechas a mano con posición, coordenadas de textura y normal (8 valores por vértice) \\ \hline
    Model & Carga modelos \texttt{.obj} con Assimp y sus texturas; también dibuja el modelo con otra textura, para la ropa de los personajes \\ \hline
    Texture & Carga imágenes con stb\_image (\texttt{LoadTexture} para jpg y \texttt{LoadTextureA} para png con transparencia) \\ \hline
    Shader & Compila los shaders y guarda las ubicaciones de los uniforms \\ \hline
    Light, DirectionalLight, PointLight, SpotLight & Las luces del modelo de Phong; \texttt{SpotLight} hereda de \texttt{PointLight} \\ \hline
    Material & Intensidad especular y brillo de cada superficie \\ \hline
    Skybox & Cubo con las seis caras del cielo \\ \hline
    Sonido & Un audio con miniaudio: tocar, pausar, repetir en bucle y, si tiene posición, sonar más bajo al alejarse la cámara \\ \hline
  \end{tabular}
  \caption{Clases del proyecto.}
\end{table}

\subsection{Funcionamiento del programa}

\begin{enumerate}
  \item \textbf{Inicio.} Se crea la ventana, se compilan los shaders y se arman las mallas hechas a mano (pisos, fachadas, jardineras y papel picado).
  \item \textbf{Carga.} Se cargan las texturas, los 84 modelos, los tres cielos y los audios, y se colocan las luces. Esta es la parte más tardada, de unos 30 segundos.
  \item \textbf{Ciclo principal}, que se repite en cada cuadro:
  \begin{enumerate}
    \item Se calcula cuánto tiempo pasó desde el cuadro anterior, para que todo se mueva igual de rápido sin importar la velocidad de la computadora.
    \item Se leen el teclado y el mouse y se mueve la cámara activa.
    \item Si el desfile no está en pausa, avanza por el recorrido, se actualizan sus KeyFrames (las chinas y la marmota) y se mueven los dulces que van en el aire. Luego se pone la cámara del desfile junto al primero de la fila.
    \item Se busca la carta más cercana a la cámara que todavía no se haya tomado (a menos de 8 metros) y se revisan las teclas \tecla{I} y \tecla{L} para tomarla, voltearla, guardarla o ver el tablero.
    \item El oído del audio se pone en la cámara activa, para que la banda y los efectos se escuchen según dónde está y hacia dónde mira.
    \item Según la hora se escoge el cielo, el color de la luz del sol o de la luna y cuántas farolas se prenden. En el modo Día de Muertos se usan su propio cielo y sus luces, y las velas parpadean.
    \item Se dibuja primero el cielo, luego los objetos opacos (entre ellos los personajes, pieza por pieza) y al final los que tienen transparencia (árboles, papel picado, telones, puestos y cartas).
    \item Al último, encima de todo, se dibuja lo que va en la pantalla, como el aviso para tomar una carta, la carta abierta o el tablero. Para esto se usa una proyección ortográfica y no se toma en cuenta la profundidad.
    \item Se muestra la imagen en la ventana.
  \end{enumerate}
\end{enumerate}

% =====================================================================
\section{Tecnologías utilizadas}

\begin{table}[H]
  \centering
  \small
  \begin{tabular}{|p{0.24\linewidth}|p{0.68\linewidth}|}
    \hline
    \textbf{Tecnología} & \textbf{Para qué se usó y por qué} \\ \hline
    C++20 y OpenGL 4.3 (GLSL 330) & Son el lenguaje y la librería de gráficos que se usaron en la materia, y permiten controlar cómo se dibuja cada cosa \\ \hline
    GLFW & Crear la ventana y leer el teclado y el mouse \\ \hline
    GLEW & Cargar las funciones de OpenGL \\ \hline
    GLM & Vectores y matrices para las transformaciones y la cámara \\ \hline
    Assimp & Cargar los modelos \texttt{.obj} con sus materiales \\ \hline
    stb\_image & Leer las imágenes jpg y png de las texturas \\ \hline
    miniaudio & Reproduce la música, los efectos y la narración, y hace que suenen según dónde está la cámara. Se eligió porque es un solo archivo y no hay que instalar nada más \\ \hline
    Visual Studio 2026 & Compilador y depurador. Es el que se usa en el curso y con el que se entrega el proyecto \\ \hline
    VS Code y CMake & Se usaron en la última parte del desarrollo para editar y compilar, porque son más ligeros. El proyecto de Visual Studio se dejó igual \\ \hline
    GIMP & Ajustar las texturas (color, recortes, orillas y transparencias). Fue el editor de imágenes principal, como se había planeado en la propuesta \\ \hline
    Blender 5.2 & Revisar los modelos de internet, pasarlos a metros, quitarles triángulos y convertirlos a \texttt{.obj}, y también armar los edificios a partir de fotos \\ \hline
    Git y GitHub & Control de versiones. Cada fase terminada se guardó en un commit \\ \hline
    LaTeX & Escribir el reporte de avances y estos manuales \\ \hline
  \end{tabular}
  \caption{Tecnologías del proyecto.}
\end{table}

\subsection{Comparación con lo planeado}

En la propuesta se compararon tres opciones para cada parte: Blender, Maya y 3ds Max para modelar, OpenGL, Unity y Unreal para programar la escena, y GIMP, Photoshop y Krita para las texturas. Al final se usaron las mismas que se habían elegido, que eran Blender, OpenGL con C++ y GIMP. Blender, porque es gratuito y exporta archivos \texttt{.obj} con normales y coordenadas de textura, que es lo que lee Assimp. OpenGL, porque es lo que se vio en la materia y obliga a entender cómo funciona cada parte (luces, cámaras y transformaciones) en vez de dejarle todo a un motor como Unity.

Lo único que se agregó fueron VS Code y CMake. Como el proyecto está en Google Drive, Visual Studio a veces fallaba al escribir sus archivos temporales en la carpeta sincronizada, y con CMake todos los archivos de la compilación quedan fuera de Drive.



% =====================================================================
\section{Técnicas de computación gráfica}

\subsection{Modelado geométrico}

Las cosas planas o que se repiten mucho se hicieron a mano con vértices en \texttt{main.cpp}. Así se armaron los pisos del andador, las banquetas y las plazas, las fachadas de las casas, las jardineras de cantera con su tierra y las tiras de papel picado. Cada tira se divide en 12 tramos para que pueda colgar más bajo en medio, como un cordel de verdad.

En los pisos, las coordenadas de textura se sacan de la posición de cada vértice en el mundo (X y Z) y no de las esquinas del rectángulo. De esta forma las losas se ven del mismo tamaño en todos los pisos, aunque cada uno mida distinto, y no se nota dónde termina uno y empieza el otro.

Para las fachadas se usan pocas imágenes en tonos claros que se pintan con el color de cada casa. Con eso salen muchas casas diferentes sin tener que cargar una textura para cada una.

\subsection{Modelos importados}

Los edificios principales (Santo Domingo, la Sangre de Cristo, el kiosco y la fuente) se armaron en Blender a partir de fotos del lugar. Los árboles, las plantas, las farolas y algunas artesanías se descargaron de Sketchfab, Poly Haven y ambientCG, todos con licencia CC0 o CC-BY (sección \ref{sec:creditos}). Antes de usarlos se pasaron a metros, a los más pesados se les quitaron triángulos y sus texturas se juntaron en menos imágenes para que cargaran más rápido.

Los modelos que se repiten, como las bancas, las macetas, los árboles y las farolas, se cargan una sola vez y se dibujan las veces que hagan falta con la función \texttt{DibujarModelo}, a la que solo se le cambian la posición, el giro y la escala.

\subsection{Modelado jerárquico}

En los personajes, cada parte del cuerpo se dibuja partiendo de la matriz de la parte de la que cuelga. El torso es la raíz y de él salen la cabeza, los brazos y las piernas, y de cada brazo sale su antebrazo, de modo que al girar el hombro el antebrazo se mueve junto con él. Los ángulos de todas las articulaciones se guardan en una estructura (\texttt{PosePersona}), y así las animaciones y los KeyFrames solo tienen que cambiar esos valores.

Los personajes parten de dos modelos de personas ya hechos (sección \ref{sec:creditos}), a los que se les hicieron muchos cambios en Blender. Se les quitaron triángulos, se partieron en piezas y a cada pieza se le puso el origen en la articulación donde gira, para que en el programa baste con trasladarla a ese punto y rotarla.

Al partir el modelo, las piezas quedaron abiertas en los cortes, y cuando un brazo se levantaba o se doblaba se veían huecos en el hombro y en el codo. Para taparlos se le puso al brazo una bola con la tela de la manga y al antebrazo una con la piel, cada una centrada en el punto donde gira su pieza. Como ese punto no se mueve cuando la pieza gira, la bola siempre queda tapando la unión. Además, se cerraron las orillas abiertas del torso y de los brazos.

\begin{itemize}
  \item \textbf{Hombre de manta y sombrero}, partido en 10 piezas: el torso (que es la raíz), la cabeza, y de cada lado el brazo, el antebrazo, el muslo y la pierna. Con él se hacen todos los hombres de la escena, que son el danzante de la marmota, el que carga el torito, los músicos y los espectadores. Entre uno y otro solo cambian la pose de los brazos y lo que llevan cargando.
  \item \textbf{China oaxaqueña}, armada por piezas igual que el hombre. Tiene el torso (con la cabeza y la blusa), la falda, la canasta y los mismos brazos del hombre, pero a 0.9 de su tamaño.
  \item \textbf{Espectadores}, que son gente común y no llevan traje de fiesta. El hombre usa las mismas piezas que el de manta, pero con un torso sin paliacate, camisa, pantalón y zapatos. La mujer usa el torso de la china con otra blusa, los brazos y las pantorrillas del hombre a 0.9 de su tamaño y una falda lisa.
\end{itemize}

Para que no todos se vean iguales, una misma pieza se puede dibujar con distintas texturas. Para eso se le agregó a la clase \texttt{Model} un segundo \texttt{RenderModel} que recibe una textura y la usa en todas las mallas, en lugar de la que trae el modelo. Así hay cuatro juegos de ropa para el hombre (camisa azul, vino, verde o blanca con pantalón de mezclilla o caqui), cuatro para la mujer (blusa blanca, rosa, amarilla o turquesa), y cada china lleva la falda de un color distinto (magenta, roja, verde y azul). La falda de la mujer es blanca y se pinta con el color del objeto, igual que los globos. Las texturas de ropa son la misma imagen original del modelo con la tela pintada de otro color, por eso las coordenadas de textura siguen funcionando.

Lo que llevan los personajes (la marmota y su palo, el torito con sus dos ruedas de cohetes y sus chispas, la trompeta, el dulce y los globos) se dibuja con la matriz de la parte del cuerpo que lo sostiene, y así se mueve junto con ella.

\begin{figure}[H]
  \centering
  \small
  \begin{tabular}{|p{0.44\linewidth}|p{0.44\linewidth}|}
    \hline
    \textbf{Hombre (\texttt{DibujarHombre})} & \textbf{Lo que se le cuelga} \\ \hline
    \parbox[t]{\linewidth}{\ttfamily
      base (posición en el recorrido)\\
      \hspace*{1em}torso\\
      \hspace*{2em}cabeza\\
      \hspace*{2em}brazo izq. \(\rightarrow\) antebrazo izq.\\
      \hspace*{2em}brazo der. \(\rightarrow\) antebrazo der.\\
      \hspace*{2em}muslo izq. \(\rightarrow\) pierna izq.\\
      \hspace*{2em}muslo der. \(\rightarrow\) pierna der.\vspace{2pt}} &
    \parbox[t]{\linewidth}{\ttfamily
      torso \(\rightarrow\) palo \(\rightarrow\) marmota\\
      torso \(\rightarrow\) torito \(\rightarrow\) 2 ruedas\\
      \hspace*{1em}\(\rightarrow\) 4 chispas c/u\\
      cabeza \(\rightarrow\) trompeta\\
      base \(\rightarrow\) 3 globos\vspace{2pt}} \\ \hline
    \textbf{China (\texttt{DibujarChina})} & \\ \hline
    \parbox[t]{\linewidth}{\ttfamily
      base\\
      \hspace*{1em}torso\\
      \hspace*{2em}falda\\
      \hspace*{2em}canasta\\
      \hspace*{2em}brazo izq. \(\rightarrow\) antebrazo izq.\\
      \hspace*{2em}brazo der. \(\rightarrow\) antebrazo der.\vspace{2pt}} &
    \parbox[t]{\linewidth}{\ttfamily
      base \(\rightarrow\) dulce (mientras vuela)\vspace{2pt}} \\ \hline
    \textbf{Espectadora (\texttt{DibujarEspectadora})} & \\ \hline
    \parbox[t]{\linewidth}{\ttfamily
      base\\
      \hspace*{1em}pantorrilla izq. y der.\\
      \hspace*{1em}cintura\\
      \hspace*{2em}falda\\
      \hspace*{2em}torso (voltea)\\
      \hspace*{3em}brazo izq. \(\rightarrow\) antebrazo izq.\\
      \hspace*{3em}brazo der. \(\rightarrow\) antebrazo der.\vspace{2pt}} &
    \parbox[t]{\linewidth}{\ttfamily
      base \(\rightarrow\) 3 globos\vspace{2pt}} \\ \hline
  \end{tabular}
  \caption{Jerarquía de los personajes: cada pieza se dibuja con la matriz de la de arriba.}
\end{figure}

\dosfotos{final/04_marmota}{Danzante con la marmota}{img/musicos_torito}{Músicos y, detrás, el torito}{El mismo hombre por piezas con distintas poses y objetos.}

\subsection{Texturizado}

\begin{itemize}
  \item Las texturas usan \textit{mipmaps} (\texttt{GL\_LINEAR\_MIPMAP\_LINEAR}), que son versiones más chicas de la imagen que OpenGL usa cuando el objeto está lejos. Sin ellos, los pisos y las fachadas parpadean cuando se mueve la cámara.
  \item Se repiten con \texttt{GL\_REPEAT}, y así una imagen chica alcanza para cubrir un piso grande.
  \item Las imágenes png con transparencia (hojas, papel picado, rejas y telones) se recortan en el shader. Si el alfa de un pixel es menor a 0.1 se descarta con \texttt{discard}, y así no tapa lo que está detrás, aunque eso se dibuje después.
  \item Los jpg deben tener un ancho múltiplo de 4, porque OpenGL lee cada fila en bloques de 4 bytes. Si no, el programa lee fuera de la imagen y se cierra.
\end{itemize}

\subsection{Iluminación}

La iluminación usa el modelo de Phong (luz ambiental, difusa y especular) en el shader \texttt{shader\_light.frag}, con tres tipos de luz:

\begin{itemize}
  \item \textbf{Una luz direccional}, que en la tarde es el sol, naranja y bajo desde el poniente para que le pegue de frente a Santo Domingo, y en la noche es la luna, azul y más débil.
  \item \textbf{16 luces puntuales}, una en cada farola (12 en el andador y 4 en el Zócalo), que solo se prenden de noche.
  \item \textbf{Tres reflectores} (\textit{spotlights}), uno que alumbra la fachada de Santo Domingo, otro para la Sangre de Cristo y otro debajo del techo del kiosco.
\end{itemize}

Al presionar la tecla N, dentro del ciclo principal se cambian el cielo, el color y la fuerza del sol o de la luna, y cuántas farolas están prendidas.

\dosfotos{retoque_texturas/02_santo_domingo}{Tarde}{retoque_texturas/08_santo_domingo_noche}{Noche}{La iluminación de la tarde y de la noche en Santo Domingo.}

\subsection{Cielo, transparencias y cámaras}

El cielo es un \textit{skybox}, que es un cubo con una imagen en cada cara (en el orden derecha, izquierda, abajo, arriba, atrás y adelante) que se dibuja primero y siempre alrededor de la cámara. Hay tres, uno de tarde, uno de noche y uno para el modo Día de Muertos, y sus caras son de 2048 pixeles.

Para las transparencias se activa el \textit{blending} (\texttt{GL\_SRC\_ALPHA}, \texttt{GL\_ONE\_MINUS\_SRC\_ALPHA}) y los objetos transparentes se dibujan al final, después de todos los opacos.

Cada cámara fija es un objeto \texttt{Camera} con su propia posición y dirección. Cuando se suelta una con la tecla F, la cámara libre copia su posición y hacia dónde mira, y se puede seguir caminando desde ahí.

% =====================================================================
\section{Animaciones}
\label{sec:animaciones}

En la propuesta se plantearon cinco animaciones para la parte de teoría, todas sobre cosas que pasan en una calenda. Aparte de esas, la china oaxaqueña y la marmota se animan con KeyFrames, y las cartas de lotería también tienen su propia animación.

\begin{table}[H]
  \centering
  \small
  \begin{tabular}{|p{0.04\linewidth}|p{0.27\linewidth}|p{0.42\linewidth}|p{0.15\linewidth}|}
    \hline
    \textbf{\#} & \textbf{Animación} & \textbf{Qué hace y para qué sirve} & \textbf{Técnica} \\ \hline
    1 & Caminata del desfile & Los personajes avanzan del atrio al Zócalo moviendo las piernas, se detienen frente al kiosco y el ciclo vuelve a empezar & Jerárquica y por estados \\ \hline
    2 & Torito & El armazón se mece sobre la cabeza del danzante y sus ruedas de cohetes giran echando chispas & Jerárquica \\ \hline
    3 & Globos & Los globos suben, se mecen y bajan sobre el desfile & Transforma\-ciones \\ \hline
    4 & Músico & Los brazos del músico suben y bajan el instrumento al tocar & Jerárquica \\ \hline
    5 & Espectadores & Los espectadores a los lados saludan al paso de la calenda & Jerárquica \\ \hline
    -- & China oaxaqueña & Saca dulces de su canasta y los lanza a la gente mientras la falda se mece & KeyFrames \\ \hline
    -- & Marmota & Gira sobre su palo mientras el palo se inclina a los lados & KeyFrames \\ \hline
    -- & Cartas de lotería & Flotan, giran y suben y bajan en su lugar para llamar la atención; al tomarlas crecen en la pantalla y dan media vuelta para mostrar la ficha & Transforma\-ciones \\ \hline
  \end{tabular}
  \caption{Animaciones del proyecto.}
\end{table}

Todas las animaciones se calculan con el tiempo real que pasa entre un cuadro y otro y no con el número de cuadros, para que se vean a la misma velocidad en cualquier computadora. El desfile tiene su propio reloj (\texttt{tiempoDesfile}), que solo avanza cuando no está en pausa. Con la tecla \tecla{P} se congelan todos los personajes del desfile y con \tecla{R} el desfile regresa al inicio. Los espectadores usan el reloj normal, así que siguen saludando aunque el desfile esté en pausa.

\subsection{Caminata y recorrido}

El recorrido se marcó con cuatro puntos en el piso, unidos con rectas: la puerta de Santo Domingo, la esquina del atrio con el andador, la entrada al andador y el frente del kiosco. La variable \texttt{avanceDesfile} guarda cuántos metros lleva el primero de la fila y aumenta 2 metros por segundo. Con la función \texttt{PuntoRecorrido} esa distancia se convierte en un punto sobre las rectas, y \texttt{GiroRecorrido} calcula hacia dónde debe voltear cada personaje usando un punto medio metro atrás y otro medio metro adelante. Gracias a eso los personajes dan la vuelta poco a poco en las esquinas, en lugar de girar de golpe.

De cada personaje se guarda en arreglos qué tipo es, cuántos metros va detrás del primero y qué tanto se hace a un lado, y con eso se dibuja en el lugar que le toca. El desfile tiene dos estados. Está \textbf{caminando} mientras el primero no llegue al kiosco, y después se queda \textbf{detenido} frente al kiosco durante 8 segundos. Cuando termina la espera, el avance regresa a cero y el desfile vuelve a salir del atrio.

Para caminar, los muslos se mueven siguiendo una curva de seno, uno hacia adelante y el otro hacia atrás. Las rodillas solo se doblan cuando la pierna va regresando, y la cintura sube y baja un poco. El ángulo de las piernas se calcula con la distancia que ya se avanzó y no con el tiempo, para que los pies no se vean resbalando sobre el piso. Los brazos se quedan quietos porque todos van cargando algo.

\subsection{Torito, músicos, globos y espectadores}

Estas animaciones se hacen con senos y cosenos que dependen del tiempo, y se aplican a las piezas del modelo jerárquico:

\begin{itemize}
  \item \textbf{Torito}: el armazón gira 8 grados a cada lado y se inclina 4 grados hacia adelante y hacia atrás. Cada rueda de cohetes da una vuelta por segundo, en sentido contrario a la otra, y sus cuatro chispas cambian de tamaño y de color como si fueran saltando.
  \item \textbf{Músicos}: cada 2 segundos levantan la cabeza y la trompeta. Hay dos poses de los brazos, una con la trompeta al frente y otra con la trompeta arriba, y la función \texttt{MezclarBrazos} pasa de una a otra poco a poco. Cada músico empieza en un momento distinto para que no se muevan todos igual.
  \item \textbf{Globos}: cada globo sube 2 metros y vuelve a bajar en 6 segundos, mientras se mueve un poco de lado y se inclina. Los tres globos de cada espectador van desfasados para que no suban al mismo tiempo.
  \item \textbf{Espectadores}: con el brazo arriba mueven la mano de un lado a otro, y cada dos saludos voltean la cabeza hacia el otro lado. Como la cabeza de la mujer viene pegada al torso, ella voltea todo el cuerpo de la cintura para arriba.
\end{itemize}

\dosfotos{img/musicos_torito}{Músicos tocando y el torito}{img/espectadores_globos}{Espectador con globos}{Animaciones de los músicos, el torito y los espectadores.}

\unafoto{img/espectadores_ropa}{Chinas con faldas de distintos colores y espectadores con ropa de diario: a la derecha una mujer de blusa amarilla y a la izquierda un hombre de camisa azul con globos.}

\subsection{KeyFrames}

Los KeyFrames se hicieron igual que en el laboratorio. Cada cuadro clave guarda sus valores, la función \texttt{interpolation} calcula cuánto tiene que cambiar cada valor en cada paso para llegar al siguiente cuadro clave, y \texttt{animate} va sumando ese cambio paso a paso. Cuando se llega al último cuadro clave se regresa al primero, por eso el último se dejó igual al primero y la animación se repite sin saltos. La diferencia con el laboratorio es que aquí se necesitaban varias animaciones al mismo tiempo, así que sus variables (\texttt{KeyFrame}, \texttt{FrameIndex}, \texttt{playIndex}, \texttt{i\_curr\_steps} e \texttt{i\_max\_steps}) se juntaron en un \texttt{struct AnimacionKeyFrames}. Los cuadros clave no se leen de un archivo, sino que se guardan en el código con \texttt{GuardarCuadro} antes del ciclo principal, y cada paso dura 1/60 de segundo.

\begin{table}[H]
  \centering
  \small
  \begin{tabular}{|p{0.15\linewidth}|p{0.27\linewidth}|p{0.47\linewidth}|}
    \hline
    \textbf{Animación} & \textbf{Qué guarda cada cuadro} & \textbf{Cuadros} \\ \hline
    China (4) & Ángulos del brazo derecho: hacia adelante, hacia el lado y del codo & 5 cuadros con 36 pasos entre cada uno (2.4 segundos por vuelta): brazo abajo, mano en la canasta, lanza al frente, abajo y descansa. Cada china empieza en otro cuadro \\ \hline
    Marmota & Giro sobre el palo e inclinación del palo & 5 cuadros con 60 pasos entre cada uno (4 segundos por vuelta): un cuarto de vuelta en cada tramo mientras el palo se inclina 4 grados a un lado y al otro \\ \hline
  \end{tabular}
  \caption{Animaciones con KeyFrames.}
\end{table}

El dulce sale volando justo cuando la china pasa del cuadro en que lo toma de la canasta al cuadro en que lo lanza. Vuela durante un segundo siguiendo una parábola hacia el frente y un poco a un lado, y va dando vueltas.

\dosfotos{final/03_chinas_musicos}{Chinas lanzando dulces}{final/12_camara_desfile}{Marmota girando}{Las dos animaciones con KeyFrames.}

% =====================================================================
\section{Modo Día de Muertos}
\label{sec:modomuertos}

La tecla \tecla{O} cambia la variable \texttt{modoMuertos}, y todo lo de este modo se dibuja dentro de un \texttt{if (modoMuertos)} en el ciclo principal. No se hizo una escena aparte. El desfile, las cámaras, las cartas y el sonido siguen funcionando igual, y lo único que cambia es que se agregan adornos y se usan otro cielo, otras luces y algunas texturas distintas. El porqué de cada adorno se explica en la sección \ref{sec:cultura}.

\dosfotos{muertos/02_velas_flotantes}{Velas flotantes y cortinas de cempasúchil}{muertos/11_altar}{Altar junto al kiosco}{El modo Día de Muertos.}

\subsection{Iluminación}

Esta noche es más oscura que la normal. El cielo es morado con un resplandor naranja y la luz direccional está casi apagada, así que casi toda la luz de la escena viene de las velas y de focos de colores.

\begin{itemize}
  \item \textbf{Farolas}: son las mismas 16, pero con una luz más débil y más naranja.
  \item \textbf{19 luces puntuales de velas}, una por cada grupo de velas (sobre las velas del andador, en las cortinas, en el kiosco, en el altar y una en cada ofrenda). Para que parpadeen como una flama, en cada cuadro su intensidad se multiplica por dos senos de distinta velocidad, y cada luz empieza en un momento distinto.
  \item \textbf{6 reflectores de colores}: Santo Domingo en morado, el kiosco en naranja, la Sangre de Cristo en turquesa, dos sobre el altar y uno sobre la catrina.
\end{itemize}

Con esto llega a haber 35 luces puntuales y 6 reflectores al mismo tiempo, por lo que los límites \texttt{MAX\_POINT\_LIGHTS} y \texttt{MAX\_SPOT\_LIGHTS} se subieron a 40 y 8, en \texttt{CommonValues.h} y también en el shader.

\subsection{Cómo se dibuja sin que el programa se vuelva lento}

\begin{itemize}
  \item \textbf{Racimos en un solo modelo.} Las velas, flores y adornos que van juntos (las velas de un cable, una cortina de flores o una ofrenda) están unidos en un solo \texttt{.obj}, y así cada racimo se dibuja con una sola llamada en lugar de una por cada vela o flor.
  \item \textbf{Una sola textura.} Todos los adornos de este modo usan la misma imagen (\texttt{muertos.png}), que tiene en distintas partes los colores de la cera, las flores, la madera y el papel. Así no hay que cambiar de textura entre un adorno y otro.
  \item \textbf{Velas que no reciben luz.} Las velas encendidas y sus flamas están en un modelo aparte de cada racimo y se dibujan al final, en una segunda vuelta, con el uniform \texttt{sinLuz} en 1. Con eso el shader usa el color de la textura tal cual, sin hacer el cálculo de Phong, y las flamas se ven brillantes aunque no les llegue ninguna luz, porque una flama de verdad brilla por sí misma.
  \item \textbf{Mismas piezas, otra textura.} Las caras de calavera de los espectadores y del desfile son las mismas texturas de ropa con la cara pintada, y se usan con el segundo \texttt{RenderModel} de la clase \texttt{Model}, igual que la ropa de colores.
\end{itemize}

\subsection{Altar y tapete de arena}

El altar tiene siete escalones con una greca al frente, y en cada escalón hay macetas de cempasúchil y de mano de león, velas, calaveritas de azúcar, pan y naranjas. Las macetas son de barro, con una mata de hojas y flores encima, y las calaveritas tienen forma de cráneo con la cara pintada al frente.

Frente al altar hay un tapete de arena de 12 metros, como los que se hacen en las calles de Oaxaca en estas fechas. Está formado por tres tramos de 3 por 4 metros con el mismo dibujo, que tiene un sol con rayos, una flor de cempasúchil grande entre flores blancas y hojas, y una greca entre filas de cuadros. A los lados tiene macetas y flores sueltas, y dos xoloitzcuintles miran hacia el altar desde distintos lugares.

\unafoto{muertos/10_tapete}{Tapete de arena frente al altar, con las macetas de cempasúchil y los xoloitzcuintles.}

\dosfotos{final/11_kiosco_muertos}{Kiosco con cempasúchil}{muertos/08_desfile_catrinas}{Desfile con caras de calavera}{Adornos del kiosco y el desfile en el modo Día de Muertos.}

% =====================================================================
\section{Cronograma: planeado contra real}

En la propuesta el trabajo se repartió en siete semanas, del 6 de octubre al 23 de noviembre, pero en la práctica el orden fue otro. Primero se armaron el escenario y la iluminación, después se rediseñó la escena, y los personajes y las animaciones quedaron para el final.

\begin{figure}[H]
  \centering
  \footnotesize
  \begin{ganttchart}[
      time slot format=isodate,
      x unit=0.205cm, y unit chart=0.5cm, y unit title=0.55cm,
      vgrid={*6{draw=none}, *1{linea, dotted}},
      title/.style={draw=linea, fill=fondo},
      title label font=\color{texto}\footnotesize,
      group/.style={draw=morado, fill=morado},
      group label font=\color{morado}\bfseries\footnotesize,
      bar/.style={draw=azul, fill=azul},
      bar label font=\color{texto}\footnotesize,
      bar incomplete/.append style={fill=fondo, draw=rosa},
      progress label text={},
      milestone/.style={draw=rosa, fill=rosa},
      milestone label font=\color{rosa}\footnotesize,
      group right shift=0, group top shift=0.6, group height=0.15
    ]{2026-10-06}{2026-11-23}
    \gantttitle{Octubre}{26} \gantttitle{Noviembre}{23} \\
    \ganttgroup{Planeado (propuesta)}{2026-10-06}{2026-11-23} \\
    \ganttbar{Investigación y estructura}{2026-10-06}{2026-10-12} \\
    \ganttbar{Santo Domingo y fachadas}{2026-10-13}{2026-10-19} \\
    \ganttbar{Personajes}{2026-10-20}{2026-10-26} \\
    \ganttbar{Integración, texturas y luces}{2026-10-27}{2026-11-02} \\
    \ganttbar{KeyFrames}{2026-11-03}{2026-11-09} \\
    \ganttbar{Cámaras, audio y pruebas}{2026-11-10}{2026-11-16} \\
    \ganttbar{Documentación}{2026-11-17}{2026-11-23} \\
    \ganttgroup{Real}{2026-10-06}{2026-11-18} \\
    \ganttbar{Propuesta e investigación}{2026-10-06}{2026-10-13} \\
    \ganttbar{Base, escenario y luces (fases 1 y 2)}{2026-10-10}{2026-10-17} \\
    \ganttbar{Cámaras (fase 3)}{2026-10-19}{2026-10-19} \\
    \ganttbar{Personaje jerárquico (fase 4)}{2026-10-20}{2026-10-21} \\
    \ganttbar{Rediseño: andador}{2026-11-02}{2026-11-04} \\
    \ganttmilestone{Reporte de avances}{2026-11-03} \\
    \ganttbar{Edificios y Zócalo}{2026-11-03}{2026-11-05} \\
    \ganttbar{Vegetación y adornos}{2026-11-05}{2026-11-11} \\
    \ganttbar{Retoque de texturas y cielo}{2026-11-11}{2026-11-13} \\
    \ganttbar{Personajes y animaciones}{2026-11-13}{2026-11-13} \\
    \ganttbar{Audio}{2026-11-13}{2026-11-14} \\
    \ganttbar{Información cultural (cartas)}{2026-11-14}{2026-11-14} \\
    \ganttbar{Manuales}{2026-11-13}{2026-11-17} \\
    \ganttbar{Videos y pruebas finales}{2026-11-14}{2026-11-17} \\
    \ganttmilestone{Entrega}{2026-11-18}
  \end{ganttchart}
  \caption{Cronograma de la propuesta (arriba) contra el real (abajo).}
\end{figure}

\begin{table}[H]
  \centering
  \small
  \begin{tabular}{|p{0.24\linewidth}|p{0.14\linewidth}|p{0.24\linewidth}|p{0.27\linewidth}|}
    \hline
    \textbf{Actividad} & \textbf{Planeado} & \textbf{Real} & \textbf{Diferencia} \\ \hline
    Investigación y estructura & 6 al 12 oct & 6 al 17 oct & Casi una semana más, porque la base del programa se armó al mismo tiempo que la investigación \\ \hline
    Santo Domingo y fachadas & 13 al 19 oct & 2 al 5 nov & Casi tres semanas después, por la pausa y el rediseño \\ \hline
    Personajes & 20 al 26 oct & 20 y 21 oct (personaje de prueba) y 13 nov (finales) & El personaje de prueba se hizo a tiempo, pero los definitivos se dejaron para el final \\ \hline
    Integración, texturas y luces & 27 oct al 2 nov & 10 al 17 oct (luces) y 5 al 13 nov (texturas) & Las luces se adelantaron y las texturas se retocaron al final \\ \hline
    KeyFrames & 3 al 9 nov & 13 nov & Cuatro días después de lo planeado \\ \hline
    Cámaras, audio y pruebas & 10 al 16 nov & 19 oct y 5 nov (cámaras), 13 y 14 nov (audio), 14 al 17 nov (pruebas) & Las cámaras se adelantaron, y el audio y las pruebas se hicieron casi en la semana planeada \\ \hline
    Documentación & 17 al 23 nov & 13 al 17 nov & Se adelantó porque la entrega quedó el 18 de noviembre \\ \hline
  \end{tabular}
  \caption{Fechas planeadas y reales de cada actividad de la propuesta.}
\end{table}

\subsection{Diferencias y sus causas}

\begin{itemize}
  \item \textbf{El escenario se adelantó y los personajes se atrasaron.} En la propuesta los personajes iban en la tercera semana, pero se decidió tener primero la escena completa (calles, edificios, luces y cámaras), porque el recorrido del desfile y el lugar de cada personaje dependen de ella.
  \item \textbf{Rediseño de la escena.} Cuando se cambió el circuito por el andador Macedonio Alcalá, el 2 de noviembre, hubo que rehacer los pisos, las casas y la posición de casi todo, pero a cambio la escena quedó mucho más parecida al lugar real.
  \item \textbf{Los modelos y las texturas tomaron más tiempo.} Los edificios se armaron a partir de fotos, y cada modelo de internet se tuvo que revisar, ajustar y probar con la luz del programa antes de usarlo. Además, al final se retocaron las texturas para que todas tuvieran el mismo estilo.
  \item \textbf{Pausa del 22 de octubre al 1 de noviembre.} Durante esas dos semanas no se avanzó porque hubo que atender entregas de otras materias. Al regresar se aprovechó para comparar la escena con fotos del lugar, y de ahí salió la idea del rediseño.
  \item \textbf{El audio y la información cultural quedaron al final.} Se dejaron para después de los personajes porque dependen de ellos, ya que la banda suena desde uno de los músicos y cuatro de las cartas van caminando con los personajes del desfile.
  \item \textbf{La documentación se adelantó.} La propuesta dejaba los manuales para la última semana, hasta el 23 de noviembre, pero la entrega quedó el 18. Por eso se empezaron el 13 de noviembre, junto con los personajes, y se fueron completando conforme se terminaba cada parte.
\end{itemize}

% =====================================================================
\section{Roles y responsabilidades}
\label{sec:roles}

\begin{table}[H]
  \centering
  \small
  \begin{tabular}{|p{0.5\linewidth}|>{\centering\arraybackslash}p{0.2\linewidth}|}
    \hline
    \textbf{Actividad} & \textbf{Alejandro} \\ \hline
    Investigación y documentación & R / A \\ \hline
    Modelado 3D y texturizado & R / A \\ \hline
    Programación OpenGL y KeyFrames & R / A \\ \hline
    Cámaras, luces y audio & R / A \\ \hline
    Pruebas, integración y manuales & R / A \\ \hline
  \end{tabular}
  \caption{Matriz RACI (R: responsable, A: aprueba).}
\end{table}

Como el proyecto se hizo de forma individual, todas las actividades tienen al mismo responsable, que también se encarga de revisarlas y aprobarlas.

Esto también explica parte del atraso con respecto al cronograma, porque el modelado de los edificios y de los personajes se tuvo que hacer al mismo tiempo que la programación. Para que alcanzara el tiempo se usaron modelos de internet con licencia libre siempre que se pudo, y los personajes se dejaron para el final.

% =====================================================================
\section{Metodología de trabajo}

En la propuesta el trabajo estaba planeado por semanas, con una actividad después de otra, pero en la práctica se trabajó por \textbf{iteraciones chicas y con un tablero tipo Kanban}:

\begin{itemize}
  \item El trabajo se dividió en fases (escenario, iluminación, cámaras, personaje jerárquico, escena del andador, personajes, animaciones, KeyFrames, audio y cierre), y después del rediseño la escena del andador se partió en cinco iteraciones chicas.
  \item Cada tarea podía estar por hacer, en curso o probada, y solo se trabajaba en una a la vez.
  \item En cada iteración se seguían los mismos pasos: ver qué modelos y texturas hacían falta, integrarlos, compilar y probar en Visual Studio, tomar capturas, escribir su parte en las evidencias y hacer su commit. No se empezaba la siguiente hasta que la anterior estaba probada y subida.
\end{itemize}

Se empezó a trabajar con iteraciones chicas después del rediseño del 2 de noviembre, porque al hacer muchos cambios grandes a la vez era difícil probarlos, y cuando algo fallaba no se sabía bien qué parte lo había causado.

Todo esto se puede ver en el historial de commits del repositorio, donde cada fase o iteración es un commit con su fecha, en las capturas de \texttt{docs/evidencias/} y en el archivo \texttt{EVIDENCIAS.md}, que tiene una sección para cada iteración.

\dosfotos{github/02_commits}{Historial de commits: cada fase o iteración es un commit con su descripción.}{github/03_commit}{Detalle de un commit con los archivos que agregó.}{Historial del repositorio en GitHub.}

Mientras se desarrollaba el proyecto, las tareas se iban anotando en la lista de fases del plan, con una casilla para cada una que se marcaba al probarla. Para la entrega, esa lista se pasó a un tablero de GitHub Projects ligado al repositorio, con las columnas \textit{Por hacer}, \textit{En progreso} y \textit{Completado}, y cada tarea tiene la fecha en que se terminó, que es la misma de su commit.

\IfFileExists{img/tablero_github.png}{%
\unafoto{img/tablero_github}{Tablero del proyecto en GitHub Projects, con las tareas por hacer, en curso y terminadas.}}{%
\pendiente{guardar la captura del tablero de GitHub Projects como \texttt{Documentacion/img/tablero\_github.png} y volver a compilar.}}

% =====================================================================
\section{Herramientas colaborativas}

\begin{itemize}
  \item \textbf{GitHub}, donde está el repositorio del proyecto, que fue privado mientras se desarrollaba y se hizo público para la entrega (\url{https://github.com/chavez-inst/ProyectoFinalCompuGraficaTeoria}). Cada fase o iteración terminada se subió como un commit que describe lo que se agregó.
  \item \textbf{GitHub Projects}, que es el tablero tipo Kanban ligado al repositorio, donde están las tareas de cada fase, su estado y la fecha en que se terminaron.
  \item \textbf{Diagrama de Gantt}, que es el que está en este documento y sirve para comparar lo planeado con lo que pasó en realidad.
\end{itemize}

\unafoto{github/01_repositorio}{Página del repositorio en GitHub a mitad del desarrollo, cuando era privado, con las carpetas del proyecto y el último commit.}

% =====================================================================
\section{Aprendizaje autónomo de herramientas nuevas}

\subsection{Alejandro: Blender, GIMP, Visual Studio y CMake}

Al empezar la materia casi no había usado programas de modelado ni de edición de imágenes (solo GIMP, para cosas sencillas), y de Visual Studio solo sabía compilar. Entre el laboratorio y este proyecto aprendí cuatro herramientas por mi cuenta, con la documentación de cada una, videos y mucha prueba y error.

\textbf{Blender.} Al principio solo lo usaba para abrir los modelos que descargaba, pero casi ninguno servía tal cual: unos venían en centímetros, otros tenían cientos de miles de triángulos y otros buscaban sus texturas en rutas que no existían. Aprendí a escalarlos a metros y aplicar la escala, a reducir triángulos con el modificador \textit{Decimate}, a revisar que las normales apuntaran hacia afuera con la vista de orientación de caras y a exportar a \texttt{.obj} con normales y coordenadas de textura para que Assimp los leyera bien. Lo más difícil fueron los personajes: para el modelado jerárquico tuve que partir cada modelo en piezas y poner el origen de cada pieza en su articulación (hombro, codo, cadera, rodilla). Ahí vi que el origen que se deja en Blender es el punto sobre el que la pieza gira en OpenGL.

\textbf{GIMP.} Lo usé para ajustar las texturas: recortarlas, corregir el color para que todas tuvieran el mismo tono, hacer que se repitieran sin que se notara la unión y crear las transparencias de las hojas, el papel picado y los telones. Tuve que aprender a trabajar con el canal alfa, porque si el png se exporta sin él, las hojas salen con su fondo. También me tocó un error difícil de encontrar: un jpg con un ancho que no era múltiplo de 4 hacía que el programa se cerrara sin ningún mensaje; desde entonces reviso el tamaño de cada imagen al exportarla.

\textbf{Visual Studio.} Aquí aprendí a configurar un proyecto completo: rutas para las librerías y para el directorio de trabajo, copiar las DLL junto al ejecutable al compilar y usar el depurador para encontrar por qué se cerraba el programa. Como el proyecto está en Google Drive para trabajar desde dos computadoras, Visual Studio a veces fallaba al escribir sus archivos temporales en la carpeta sincronizada, y aprendí a mandar la base de datos de IntelliSense a una carpeta fuera de Drive.

\textbf{CMake con VS Code.} Casi al final, por los mismos problemas con Drive, pasé a VS Code con CMake. Con la documentación de CMake y de la extensión CMake Tools escribí un \texttt{CMakeLists.txt} y dos presets (Debug y Release) que dejan todo lo que genera la compilación fuera de Drive. Lo que más me costó fue lograr que VS Code usara el compilador y el depurador de Visual Studio. El proyecto de Visual Studio se mantuvo igual, así que se puede compilar desde cualquiera de los dos.

% =====================================================================
\section{Individual reflection (English)}

I chose the Calenda of Oaxaca because I'm originally from there and it's been a part of my personal identity since I was a kid. I started playing the saxophone when I was 9 or so, and I mainly played Calenda and Guelaguetza songs. At first I had a very different idea of the project, but when I went back for the Day of the Dead season I reworked a lot of it so it would feel like the real Oaxaca. Since I worked on the project alone, I was in charge of every part: the research, the scene, the models and textures, the OpenGL program and the documentation. I checked each part against my own experience of Oaxaca so it would feel as real as possible.

One of the most difficult parts of the project was the characters: each one had to be cut into pieces with the pivot of every joint in the right place, so that the arms and legs would rotate naturally inside the hierarchy (like we did in the lab).

In the last week I added sound and the cultural information. The band plays from one of the musicians, so it sounds louder as you get closer, and the information is shown as lotería cards that you have to find along the street; so the user has to explore the whole scene instead of just reading a menu.

I learned how much planning matters in a graphics project. Even small changes can break other parts of the program, so using git made the work much easier. I also learned to prepare models in Blender before loading them with Assimp, to edit textures in GIMP so that they all share the same style, and to build the project with CMake from Visual Studio Code.

If I started again, I would build the characters and animations earlier, because the scene took more time than expected and they were left for the last week. Even so, I am happy with the final result and with how close it looks to the real Oaxaca, especially to people who know the city.

% =====================================================================
\clearpage
\section{Conclusiones}

\subsection{Técnicas}

Solo con lo que se vio en la materia (transformaciones, modelado jerárquico, texturas, iluminación de Phong, skybox, transparencias y cámaras) se pudo armar una escena de unos 300 metros de largo con 84 modelos, 16 luces puntuales, tarde y noche, un desfile de 9 personajes con siete animaciones (dos de ellas con KeyFrames), audio en 3D y diez cartas con información. Más que usar técnicas avanzadas, lo que hizo que la escena se viera bien fue cuidar los detalles, como respetar el tamaño real de cada cosa, revisar que las normales apuntaran hacia afuera, dibujar los transparentes en el orden correcto, no usar modelos con demasiados triángulos y que todas las texturas tuvieran el mismo estilo.

El programa tarda unos 30 segundos en cargar y, compilado en Release, corre bien en una computadora con tarjeta de video dedicada. Algo que no se esperaba es que preparar los modelos y las texturas tomó más tiempo que programar, así que en un proyecto como este conviene tenerlos listos desde el principio.

\subsection{Sociales}

Con un recorrido virtual se puede conocer la calenda y el centro de Oaxaca sin tener que viajar hasta allá ni esperar a julio. Solo hace falta una computadora, y como se maneja con el teclado y el mouse, cualquiera lo puede usar siguiendo el manual de usuario. Por otro lado, al ser un trabajo individual hubo que organizarse solo, y en eso las iteraciones chicas y el tablero de tareas ayudaron mucho a llevar el control del avance.

\subsection{Culturales}

La idea era que la escena se viera como Oaxaca y no como una imagen general de México. Por eso cada edificio, planta y objeto se escogió y se colocó viendo fotos y mapas del lugar real, y los datos de las fichas se sacaron de fuentes como el INAH, la UNESCO, el gobierno de Oaxaca y el Diario Oficial. Con las cartas se juntan dos tradiciones mexicanas, la calenda y la lotería, y como hay que buscarlas, el usuario termina recorriendo toda la calle y conoce cada personaje en el lugar donde está, en vez de solo leer una lista. También se revisó que la escena correspondiera a la época de la fiesta (sección \ref{sec:cultura}), y lo único que no es de julio son las jacarandas con flores.

\subsection{Económicas}

En la propuesta el costo se calculó solo con las horas de trabajo, contando dos integrantes con 40 horas cada uno, a \$100 la hora de un programador junior, lo que daba \$8,000 en total. Al final el proyecto lo hizo una sola persona. Las horas reales se calcularon a la baja con las fechas de los commits y de las sesiones de trabajo:

\begin{table}[H]
  \centering
  \small
  \begin{tabular}{|p{0.6\linewidth}|>{\centering\arraybackslash}p{0.15\linewidth}|}
    \hline
    \textbf{Etapa} & \textbf{Horas} \\ \hline
    Propuesta e investigación & 8 \\ \hline
    Base, escenario, luces y cámaras (fases 1 a 3) & 12 \\ \hline
    Personaje de prueba (fase 4) & 4 \\ \hline
    Rediseño del andador y reporte de avances & 8 \\ \hline
    Edificios y Zócalo & 10 \\ \hline
    Vegetación, adornos y retoque de texturas & 14 \\ \hline
    Personajes y animaciones & 12 \\ \hline
    Audio y cartas de lotería & 8 \\ \hline
    Manuales, videos y pruebas & 9 \\ \hline
    \textbf{Total} & \textbf{85} \\ \hline
  \end{tabular}
  \caption{Horas de trabajo aproximadas por etapa.}
\end{table}

\begin{table}[H]
  \centering
  \small
  \begin{tabular}{|p{0.4\linewidth}|>{\centering\arraybackslash}p{0.2\linewidth}|>{\centering\arraybackslash}p{0.2\linewidth}|}
    \hline
    \textbf{Concepto} & \textbf{Estimado} & \textbf{Real} \\ \hline
    Horas de trabajo & 80 (2 × 40) & 85 (1 persona) \\ \hline
    Costo de las horas (\$100 por hora) & \$8,000 & \$8,500 \\ \hline
    Licencias de software & \$0 & \$0 \\ \hline
    Modelos, texturas, audio e imágenes & -- & \$0 \\ \hline
    Electricidad de la computadora (aprox.) & -- & \$100 \\ \hline
    \textbf{Total} & \textbf{\$8,000} & \textbf{\$8,600} \\ \hline
  \end{tabular}
  \caption{Costo estimado en la propuesta contra el real, en pesos mexicanos.}
\end{table}

El costo total quedó muy cerca de lo estimado, solo un 7.5~\% más, pero una sola persona hizo el trabajo que se había pensado para dos, así que sus horas fueron más del doble de lo planeado. No hubo gastos en software ni en recursos, porque todas las herramientas son gratuitas (Visual Studio Community, VS Code, Blender, GIMP y Git) y los recursos de internet tienen licencias libres o se usan con fines escolares.

\subsection{Ambientales}

El proyecto no usa materiales físicos, y un recorrido virtual puede reemplazar algunos viajes que se harían solo para conocer un lugar. Además, al reducir los triángulos de los modelos y el tamaño de las texturas, la tarjeta de video trabaja menos y se gasta menos energía.

% =====================================================================
\section{Créditos y licencias}
\label{sec:creditos}

Los edificios principales, las casas, el kiosco, la fuente, las bancas, los carritos, los puestos y las texturas de los pisos y las fachadas se hicieron o se ajustaron para este proyecto. Estos son los recursos de otras personas que se usan en el programa:

\begin{longtable}{|p{0.28\linewidth}|p{0.2\linewidth}|p{0.13\linewidth}|p{0.27\linewidth}|}
  \hline
  \textbf{Recurso} & \textbf{Autor y fuente} & \textbf{Licencia} & \textbf{Cambios} \\ \hline
  \endhead
  Street Lamp (farola de brazo) & bez\_glaza, Sketchfab & CC BY 4.0 & Tamaño y formato \\ \hline
  Roanoke Street Lamp (farola de globos) & loganh96, Sketchfab & CC BY 4.0 & Menos triángulos, metal negro \\ \hline
  Foto del templo de Santo Domingo de Guzmán & nan palmero, Wikimedia Commons & CC BY 2.0 & Enderezada, recortada y con color corregido \\ \hline
  Cobblestone Floor 08 (empedrado) & Poly Haven & CC0 & -- \\ \hline
  Plaster001 (aplanado) y Planks013 (madera) & ambientCG & CC0 & -- \\ \hline
  Agave Sisalana Plant & Md Mahmudur Rahman Bappy, Sketchfab & CC BY 4.0 & Menos triángulos, una sola planta \\ \hline
  Chinese flame tree (flamboyán) & PlantCatalog, Sketchfab & CC BY 4.0 & Menos triángulos, hojas en planos, flores rojas \\ \hline
  Frangipani tree (cacaloxóchitl) & PlantCatalog, Sketchfab & CC BY 4.0 & Menos triángulos, hojas ajustadas \\ \hline
  Blue jacaranda & PlantCatalog, Sketchfab & CC BY 4.0 & Rehecho como árbol ligero \\ \hline
  Trees Low Poly & Igor\_K., Sketchfab & CC BY 4.0 & Dos árboles separados \\ \hline
  Tree low poly lowpoly & 00amza, Sketchfab & CC BY 4.0 & Árboles separados, hojas ajustadas \\ \hline
  Flower Pot & empierz, Sketchfab & CC BY 4.0 & Con tierra y plantas \\ \hline
  Saintpaulia flowers low-poly & Elena FF, Sketchfab & CC BY 4.0 & Solo las flores \\ \hline
  Low poly trees, flowers and grass & Márcio Meireles (MMC Works), Sketchfab & CC BY 4.0 & Flores y pasto separados \\ \hline
  Clay mug with face & Martin Teichmann, Sketchfab & CC BY 4.0 & Formato y tamaño \\ \hline
  Oaxacan Birb (alebrije) & DrLemurpants, Sketchfab & CC BY 4.0 & Menos triángulos \\ \hline
  Perrito (alebrije) & aromeromyt, Sketchfab & CC BY 4.0 & Menos triángulos, limpieza \\ \hline
  Petate & ocelotlmx, Sketchfab & CC BY 4.0 & Formato \\ \hline
  La Calavera Catrina - Day of the Dead & Brian Trepanier, Sketchfab & CC BY 4.0 & Menos triángulos, textura en una sola imagen \\ \hline
  Xolo Jorge (xoloitzcuintle) & Jorgeguillermon, Sketchfab & CC BY 4.0 & Tamaño y textura de piel \\ \hline
\end{longtable}

El alebrije de coyote está basado en una pieza de Pablo Franco y Nelly Ramírez.
Las librerías que se usan son de código abierto: GLFW (licencia zlib), GLEW (BSD y MIT), GLM (MIT), Assimp (BSD de 3 cláusulas), stb\_image (dominio público o MIT) y miniaudio (dominio público o MIT-0).

Los personajes del desfile parten de dos modelos de personas que proporcionó un compañero, sin licencia: un hombre con traje de manta y sombrero y una bailarina con traje típico. Se usaron como base, pero se modificaron bastante. El hombre se partió en piezas con el origen en cada articulación, y la china oaxaqueña se armó con la cabeza y la blusa de la bailarina, los brazos del hombre y una falda, una canasta y un dulce que se hicieron para el proyecto. La ropa de los espectadores y las faldas de colores son las mismas texturas originales con la tela pintada de otro color. La marmota, el torito con sus ruedas, la trompeta y los globos también se hicieron para este proyecto.

\subsection{Audio}

La música y los efectos se tomaron de videos públicos de YouTube. No tienen una licencia libre, así que se usan solo con fines escolares, sin fines de lucro y citando a su autor. A todos se les ajustó el volumen para que sonaran parejo.

\begin{longtable}{|p{0.3\linewidth}|p{0.25\linewidth}|p{0.35\linewidth}|}
  \hline
  \textbf{Audio} & \textbf{Autor y fuente} & \textbf{Dónde se usa y cambios} \\ \hline
  \endhead
  Son Calenda & Banda Tierra Mojada, YouTube & Banda del desfile; sin la introducción y en mono \\ \hline
  Dios nunca muere (M. Alcalá) & Banda Filarmónica del CECAM, YouTube & Banda del desfile; en mono \\ \hline
  Ambiente parque ciudad & La SonidoTeca, YouTube & Murmullo de la gente; se repite sin corte \\ \hline
  Cuete cohete con silbido & Jerry Efectos de Sonido y Audio, YouTube & Cohete al salir el desfile \\ \hline
  Cohete coloso efecto 1 & Pirotecnia La Mecha, YouTube & Cohete al llegar al kiosco \\ \hline
  Church bell ringing effects 1.0 & Xeno's Random Edits, YouTube & Campanas de Santo Domingo; solo los primeros 18 segundos \\ \hline
  Fuegos artificiales & Esperanza Digital, YouTube & Chispas del torito; se repite sin corte \\ \hline
\end{longtable}

\subsection{Cartas de lotería}

Las imágenes de las diez cartas (El Barril, La Luna, La Dama, El Diablito, El Tambor, El Sol, El Corazón, El Catrín, La Mano y El Nopal) son las de la lotería tradicional de Don Clemente. Se usan con fines escolares y lo único que se les cambió fue redondearles las esquinas. Las fichas, las pistas del tablero, los avisos y el letrero de lotería se hicieron para el proyecto.

% =====================================================================
\section{Referencias}

\begin{itemize}[label={}, leftmargin=1.27cm, itemindent=-1.27cm]
  \item Assimp Development Team. (s.~f.). \textit{Open Asset Import Library (assimp)} [Software]. GitHub. \url{https://github.com/assimp/assimp}
  \item Blender Foundation. (2026). \textit{Blender} (versión 5.2) [Software]. \url{https://www.blender.org}
  \item de Vries, J. (2020). \textit{Learn OpenGL: Learn modern OpenGL graphics programming in a step-by-step fashion}. Kendall \& Welles. \url{https://learnopengl.com}
  \item Embajada de México ante la Santa Sede. (s.~f.). \textit{La Guelaguetza}. Secretaría de Relaciones Exteriores. \url{https://embamex.sre.gob.mx/vaticano/index.php/noticias/771-la-guelaguetza}
  \item Gobierno del Estado de Oaxaca. (2026). \textit{Guelaguetza Oaxaca 2026}. \url{https://guelaguetza.oaxaca.gob.mx/}
  \item Hughes, J. F., van Dam, A., McGuire, M., Sklar, D. F., Foley, J. D., Feiner, S. K., \& Akeley, K. (2013). \textit{Computer graphics: Principles and practice} (3.ª ed.). Addison-Wesley.
  \item Instituto Mexicano de la Propiedad Industrial. (2022, 2 de marzo). Declaración de protección de la indicación geográfica Tapetes de Teotitlán, Santa Ana y San Miguel del Valle. \textit{Diario Oficial de la Federación}. \url{https://www.gob.mx/cms/uploads/attachment/file/749823/IMPI.Declaracio_n_Tapetes_T._02.03.22.pdf}
  \item Instituto Nacional de Antropología e Historia. (s.~f.). \textit{Museo de las Culturas de Oaxaca}. \url{https://lugares.inah.gob.mx/en/node/4252}
  \item Instituto Nacional de Antropología e Historia. (2023, 15 de febrero). \textit{El INAH firma dos convenios en favor del patrimonio cultural del municipio de Oaxaca y del estado} (Boletín núm. 76). \url{https://inah.gob.mx/images/boletines/2023/76/Boletin%2076.pdf}
  \item México Desconocido. (s.~f.). \textit{Las calendas de Oaxaca, explosión de color, música y tradición que inunda las calles}. \url{https://www.mexicodesconocido.com.mx/las-calendas-de-oaxaca.html}
  \item Reiss, D. (2026). \textit{miniaudio} (versión 0.11) [Software]. GitHub. \url{https://github.com/mackron/miniaudio}
  \item Sánchez Ulloa, C. A. (2022). Un acercamiento al trabajo de los puesteros de lotería de cartones en Campeche, 1904-1916. \textit{Letras Históricas}, (26), 1--31. \url{https://doi.org/10.31836/lh.26.7347}
  \item Secretaría de Turismo de Oaxaca. (2018a). \textit{San Martín Tilcajete, alebrijes} [Ficha de producto turístico OAX004]. \url{https://www.oaxaca.gob.mx/sectur/wp-content/uploads/sites/65/2020/09/PT-OAX001-2018.pdf}
  \item Secretaría de Turismo de Oaxaca. (2018b). \textit{Teotitlán del Valle, tapetes de lana} [Ficha de producto turístico OAX011]. \url{https://www.oaxaca.gob.mx/sectur/wp-content/uploads/sites/65/2020/09/PT-OAX004-2018.pdf}
  \item UNESCO. (2008). \textit{Indigenous festivity dedicated to the dead}. Intangible Cultural Heritage. \url{https://ich.unesco.org/en/RL/indigenous-festivity-dedicated-to-the-dead-00054}
  \item UNESCO World Heritage Centre. (s.~f.). \textit{Historic Centre of Oaxaca and Archaeological Site of Monte Albán}. \url{https://whc.unesco.org/en/list/415}
  \item Wright, R. S., Lipchak, B., \& Haemel, N. (2010). \textit{OpenGL SuperBible: Comprehensive tutorial and reference} (5.ª ed.). Addison-Wesley.
\end{itemize}

\end{document}
